\section{Introduction}
\label{sec:introduction}

The Internet of Things (IoT) spans critical infrastructure, yet authentication mechanisms remain underdeveloped \cite{iotreview, ahvanooey2022modern} for resource-constrained edge devices. Widespread reliance on default credentials \cite{ibm2025} has pushed the industry toward inherence-based security (i.e., biometrics) \cite{markets2025}. However, despite this adoption, the integration landscape is severely fragmented; biometric modules ship with proprietary protocols, vendor-specific drivers, and mutually incompatible APIs \cite{yang2021biometric}.

Vendor fragmentation creates an architectural gap, particularly on resource-constrained microcontrollers. While high-level operating systems comply with standards like ISO/IEC 19784 (BioAPI) \cite{iso19784}, they lean on framework runtimes like D-Bus that require far more memory than a typical IoT end-node provides. ISO/IEC 29164 (Embedded BioAPI) \cite{iso29164} standardizes hardware interfaces at the wire-protocol level but specifies no host-side software behaviour. Developers are forced to write vendor-specific and non-portable code, which also results in inconsistent security policy enforcement.

This architectural gap is particularly evident in RTOS environments. Legacy approaches force developers to use generic sensor subsystems, forcing biometrics sensors to behave like passively read sensors, leaving enrollment state management entirely to the application layer, and cramming user identities into physics data structures. These limitations validate the need for a dedicated biometric subsystem that resolves type safety violations, stateless APIs applied to transactional workflows, and no focus on user management.

While legacy RTOS frameworks merely wrap bus transactions, to our knowledge, no existing solution offers a dedicated, kernel-level biometric abstraction. The fundamental novelty of the Biometric Abstraction for Secure Edge (BASE) lies in moving enrollment state machines (multi-step), type-safe data structures, and normalized security policies into a strict driver contract. By shifting this complexity into the kernel, BASE decouples application logic from hardware temporalities, enabling zero-code-change migrations across diverse sensor modalities. BASE is the first dedicated kernel-level biometric subsystem designed specifically for resource-constrained RTOS environments.

The specific contributions of this paper are:
\begin{itemize}
    \item \textbf{Kernel-Level Biometric Abstraction:} We propose BASE, a vendor-agnostic architecture that prevents type-safety violations and broken state management by enforcing a strict, deterministic enrollment state machine directly within the driver layer.
    \item \textbf{Strict Modality Independence:} We validated BASE across three different physical sensor families: contact-based optical fingerprint (ZFM-X0, GT-5X), and camera-based multi-modal recognition (DFRobot AI10, face and palm vein). The AI10 driver was developed by an independent contributor with zero modifications to the BASE API layer, confirming third-party implementability. Zero lines of application code changed across all three families.
    \item \textbf{Deterministic Multitasking:} BASE's API structures blocking operations around bounded timeouts. Empirical evaluation on an ARM Cortex-M4 shows BASE reduces maximum scheduling jitter from \SI{490}{\milli\second} to \SI{100}{\micro\second} and eliminates deadline misses entirely (0 vs.\ 68 out of 200) compared to a legacy blocking driver.
    \item \textbf{Industrial Standardization:} This architecture has been accepted by the Zephyr Architecture Working Group, merged into the mainline OS, and released as part of Zephyr 4.4\cite{zephyr_biometrics, zephyr_pr, zephyr_pr2, zephyr_44_notes}, where it is actively maintained by the first author.
\end{itemize}

The remainder of this paper is organized as follows: Section~\ref{sec:related_work} reviews related work. Section~\ref{sec:architecture} details the BASE architecture. Section~\ref{sec:implementation} describes the Zephyr integration. Section~\ref{sec:evaluation} presents the experimental results, and Section~\ref{sec:conclusion} concludes.
